\chapter*{Vysvetlivky}

\begin{tabular}{>{\raggedright\arraybackslash}p{4cm} >{\raggedleft\arraybackslash}p{10cm}}
    \textbf{Vzdialenosť a dostrel} & Vzdialenosť medzi susednými hráčmi je 1; vždy sa počíta kratšia cesta po stole a vyradení hráči sa nezapočítavajú. Vzdialenosť môžu meniť iba karty a schopnosti, ktoré to výslovne uvádzajú (napr. Mustang, Scope). Dostrel je vzdialenosť, na ktorú hráč dosiahne kartou alebo schopnosťou. Zbrane menia iba dostrel hráča — samy o sebe nemenia vzdialenosť medzi hráčmi. \\\\

    \textbf{Ľubovoľný hráč} & Na vzdialenosti hráča od vás nezáleží. \\\\

    \textbf{Ťahať kartu} & Otočiť vrchnú kartu balíčka na vyhodnotenie situácie, napr. pri Bareli, Väzení alebo Dynamite. Takto ťahaná karta sa po vyhodnotení odhodí. \\\\

    \textbf{Potiahnuť kartu} & Získať kartu do ruky, či už z balíčka, od iného hráča alebo z iného zdroja. \\\\

    \textbf{Karta a efekt} & Karta $\neq$ efekt. Karta môže vytvárať efekt inej karty, ale stále ide o dve rozdielne veci. Ak text zakazuje alebo obmedzuje \emph{kartu}, myslí sa tým konkrétna karta s tým názvom. Ak text hovorí o \emph{efekte}, vzťahuje sa na výsledok pôsobenia bez ohľadu na to, z akej karty alebo schopnosti vznikol. Napríklad zákaz hrať kartu Pivo sa nevzťahuje automaticky na všetky efekty Pivo a požiadavka na efekt Vedle! sa nemusí vzťahovať iba na konkrétnu kartu Vedle!. \\\\

    \textbf{Fáza 1} & Fáza, v ktorej si hráč potiahne karty z balíčka.\\\\
    \textbf{Fáza 2} & Fáza hrania kariet a používania schopností. \\\\
    \textbf{Fáza 3} & Odhodenie prebytočných kariet a ukončenie ťahu. Na konci ťahu môže mať hráč na ruke najviac toľko kariet, koľko má práve životov. \\\\
\end{tabular}
